% Generated by examples.py; do not edit.
\begingroup\small\raggedright
\setlength{\parindent}{0pt}\setlength{\parskip}{3pt}
Premise 1:\par
\textasciigrave{}\textasciigrave{}\textasciigrave{}\par
The text “\{\{a\}\}” is the canonical decimal numeral for a nonnegative integer: either 0 or a digit from 1 through 9 followed by zero or more digits from 0 through 9.\par
\textasciigrave{}\textasciigrave{}\textasciigrave{}\par
\smallskip
Premise 2:\par
\textasciigrave{}\textasciigrave{}\textasciigrave{}\par
The text “\{\{b\}\}” is the canonical decimal numeral for a nonnegative integer: either 0 or a digit from 1 through 9 followed by zero or more digits from 0 through 9.\par
\textasciigrave{}\textasciigrave{}\textasciigrave{}\par
\smallskip
Premise 3:\par
\textasciigrave{}\textasciigrave{}\textasciigrave{}\par
The text “\{\{c\}\}” is the canonical decimal numeral for a nonnegative integer: either 0 or a digit from 1 through 9 followed by zero or more digits from 0 through 9.\par
\textasciigrave{}\textasciigrave{}\textasciigrave{}\par
\smallskip
Conclusion:\par
\textasciigrave{}\textasciigrave{}\textasciigrave{}\par
The request “What do you get if you multiply the sum of \{\{a\}\} and \{\{b\}\} by \{\{c\}\}?” asks a numerical question equivalent to finding the value of the well-formed closed arithmetic expression “((\{\{a\}\} + \{\{b\}\}) * \{\{c\}\})”, whose value is a nonnegative integer. For every canonical nonnegative decimal numeral n, n correctly answers the request if and only if it denotes that value.\par
\textasciigrave{}\textasciigrave{}\textasciigrave{}\par
\endgroup
